\chapter{研究点三 实际市场结果返回检验与动态监控技术}

\section{引言}

【说明本研究点解决的问题及其与总体目标的关系。】

\section{金融风险返回检验与监控相关研究}

【评述与本研究点方法设计直接相关的工作。】

\section{基于时点冻结证据链的验证与监控方法}

\subsection{研究动机与目的}

【说明方法提出的动机、适用边界和预期改进。】

\subsection{方法总览}

【给出方法总体流程、输入输出和关键假设。】

\subsection{五年历史回放与返回检验}

【填写第一项方法设计、公式、算法或实现细节。】

\subsection{动态监控、异常诊断与证据归档}

【填写第二项方法设计、公式、算法或实现细节。】

\section{实验验证与结果分析}

\subsection{评估指标}

【定义数据划分、基准、指标、统计检验和通过门槛。】

\subsection{返回检验与统计稳健性试验}

【报告对比试验设计、结果、负向证据和分析。】

\subsection{监控、恢复与证据一致性试验}

【报告补充对比、消融、稳健性或敏感性试验。】

\section{本章小结}

【总结本研究点的主要发现、限制及其对后续章节的影响。】
